\section{Introduction}\label{sec:introduction}

Goldfeld's mean-rank conjecture predicts that the average central order
of vanishing in a quadratic-twist family is \(1/2\)
\cite[p.~113, Conjecture~(B)]{Goldfeld1979}.
The issue addressed here is the contribution of rare twists of large
analytic rank: knowing the densities of ranks zero and one does not
determine the mean without control of this contribution.

Let \(E/\Q\) be an elliptic curve. For a nonzero squarefree integer \(d\),
write \(E^{(d)}\) for its quadratic twist, and put
\[
 \mathcal D(Y)=\{d\in\Z:0<|d|\le Y,\ d\text{ squarefree}\},
 \qquad
 a(E^{(d)})=\ord_{s=1}L(E^{(d)},s).
\]
Both signs of \(d\) are included. We call \(a\) the analytic rank and
write \(r(E^{(d)})=\operatorname{rank}_{\Z}E^{(d)}(\Q)\) for the
algebraic or Mordell--Weil rank. Outside the explicitly algebraic
statements below, rank means analytic order of vanishing.

\begin{theorem}\label{thm:mean}
For every elliptic curve \(E/\Q\),
\[
 \lim_{Y\to\infty}\frac{1}{\#\mathcal D(Y)}
     \sum_{d\in\mathcal D(Y)}a(E^{(d)})=\frac12.
\]
\end{theorem}

Theorem~\ref{thm:mean} resolves positively Goldfeld's mean analytic-rank
conjecture in this signed squarefree counting convention. It applies
without restrictions on complex multiplication, rational torsion,
rational isogenies, or reduction type.

The analytic density theorem
\cite[Theorem~1.2]{OpenAIAnalyticDensityV2}
states that, for this same family and every \(E/\Q\),
\begin{equation}\label{eq:density-input}
 \frac{\#\{d\in\mathcal D(Y):a(E^{(d)})=j\}}{\#\mathcal D(Y)}
       \longrightarrow\frac12
       \qquad(j=0,1).
\end{equation}
For the analytic mean theorem, \eqref{eq:density-input} is the only input
from the companion manuscript identified in the bibliography.
A zero-density exceptional family can still contribute
to a rank-weighted average. The new assertion needed for the mean is the
following tail estimate, whose proof is independent of
\eqref{eq:density-input}.

\begin{theorem}\label{thm:tail}
For every elliptic curve \(E/\Q\), there are constants \(C_E>0\) and
an integer \(R_E\ge1\) such that, for every integer \(R\ge R_E\),
\[
 \limsup_{Y\to\infty}\frac1Y
 \sum_{\substack{d\in\mathcal D(Y)\\a(E^{(d)})>R}}a(E^{(d)})
       \le \frac{C_E}{R}.
\]
\end{theorem}

Together, \eqref{eq:density-input} and Theorem~\ref{thm:tail} show that
\[
 \sum_{\substack{d\in\mathcal D(Y)\\a(E^{(d)})\ge2}}a(E^{(d)})=o_E(Y).
\]
For each fixed sufficiently large integer \(R\), \eqref{eq:density-input}
makes the complement of ranks zero and one have density zero, so ranks
\(2\) through \(R\) contribute \(o_{E,R}(Y)\).
For the remaining ranks, use the height limsup in Theorem~\ref{thm:tail}
and then let \(R\) tend to infinity. This order of limits will be
made explicit in Section~\ref{sec:tail}.

\subsection{Algebraic-rank moments}

The same signed squarefree family has the following algebraic-rank
moment limits. This consequence uses the companion's density conclusions
and its density-one equality of analytic and algebraic ranks, together
with the exponential-moment bound of Koymans and Smith. It does not use
the analytic tail estimate of Theorem~\ref{thm:tail}.

\begin{corollary}[Fixed algebraic-rank moments]\label{cor:algebraic-moments}
Fix an elliptic curve \(E/\Q\). For every fixed real \(t\),
\[
 \lim_{Y\to\infty}\frac1{\#\mathcal D(Y)}
 \sum_{d\in\mathcal D(Y)}\exp\!\bigl(t\,r(E^{(d)})\bigr)
 =\frac{1+e^t}{2}.
\]
For every fixed positive integer \(m\),
\[
 \lim_{Y\to\infty}\frac1{\#\mathcal D(Y)}
 \sum_{d\in\mathcal D(Y)}r(E^{(d)})^m=\frac12.
\]
\end{corollary}

The proof appears in Section~\ref{sec:algebraic}.

All parameters \(t\) and moment orders \(m\) in
Corollary~\ref{cor:algebraic-moments} are fixed before \(Y\to\infty\).
The corollary concerns algebraic ranks only: it gives no higher
analytic-rank moments, no uniformity for parameters varying with \(Y\),
and no moment limit for thin polynomial subfamilies.

\subsection{History and significance}

Goldfeld's original question concerns the change of rank when a fixed
elliptic curve is viewed over quadratic fields. Under
Birch--Swinnerton-Dyer, that change is the central order of the associated
quadratic twist. His Conjecture~(B) formulates the average directly in
terms of this analytic order, using quadratic-field discriminants
ordered by absolute value \cite[pp.~113--114]{Goldfeld1979}.
Here the counting parameter is instead the signed squarefree integer
\(d\) in \(\mathcal D(Y)\).
The rank-zero/rank-one density prediction describes the smallest orders
allowed by the two functional-equation signs. The mean prediction also
requires the larger orders, however sparse, to have negligible total
contribution.

Conditional analytic estimates illustrate both aspects of the problem.
Heath-Brown proved an upper bound \(3/2+o(1)\) for smooth averages over
fundamental discriminants coprime to the conductor, separately for each
functional-equation sign \cite[Theorem~3]{HeathBrown2003}.
His nonnegative compactly supported weight selects one sign of the
discriminant, and his Riemann-hypothesis assumption concerns every
quadratic twist, including those not coprime to the conductor.
Under the same twistwise Riemann-hypothesis assumption, Miller and Wong
bounded higher analytic-rank moments and obtained exponentially
decreasing large-rank counting tails
\cite[Theorem~1.1 and Corollary~1.2]{MillerWong2012}.
Their weighted sums run over integer parameters, with possible repeated
squareclasses; the weight selects one parameter sign without fixing the
functional-equation sign. Their tail estimate fixes the rank threshold
before taking the height limit. Fiorilli obtained the exact mean \(1/2\) over
signed squarefree parameters coprime to the conductor, assuming the
Riemann hypothesis for elliptic-curve \(L\)-functions and an additional
averaged cancellation hypothesis for nonreal zeros
\cite[Theorem~1.1 and Hypothesis~M]{Fiorilli2016}.
Thus this exact conditional mean uses more than the Riemann hypothesis.

A different route to densities comes from Selmer groups.
Smith proved that, for every rational elliptic curve, the full
\(2\)-power Selmer corank is zero or one with density \(1/2\) each,
using signed integer twist parameters
\cite[Theorem~1.1]{Smith2025}. His analytic density corollary assumes
Birch--Swinnerton-Dyer \cite[Corollary~1.2]{Smith2025}.
The companion's \(2\)-converse supplies the analytic implication needed
for \eqref{eq:density-input} in the signed squarefree convention.
Theorem~\ref{thm:tail} addresses the additional rank-mass question
independently of that density argument.

Koymans and Smith prove exponential moment bounds for Mordell--Weil rank
in polynomial quadratic-twist families
\cite[Theorem~1.4]{KoymansSmith2026}. The specialization in the proof of
Corollary~\ref{cor:algebraic-moments} supplies the needed algebraic
tail control. Such bounds do not supply an analytic-rank identity or
higher analytic-rank moments.

Hanners claims the full Birch--Swinnerton-Dyer conjecture for every
rational elliptic curve \cite[Theorems~28.1, 39.4 and~39.5]{Hanners2026}. We do not use
this claim. The unresolved point is how the bridge conditions supported
by tests on eighteen curves in Section~39.1 are established for every
rational elliptic curve, as required by the transfer in Section~39.2.
This is not a refutation of the claimed theorem.

Modularity provides the analytic continuation and functional equation
for every curve in the argument
\cite[Theorem~A]{BreuilConradDiamondTaylor2001}. Our estimates use these analytic properties
and the local coefficient bounds. The proof invokes neither the
Birch--Swinnerton-Dyer conjecture nor the generalized Riemann hypothesis.

Earlier work on a fixed elliptic curve already connected derivative
moments to nonvanishing and analytic-rank averages. Perelli and
Pomyka\l a proved first-derivative nonvanishing results and bounds for
the sum of analytic ranks \cite{PerelliPomykala1997};
Pomyka\l a extended this nonvanishing approach to higher derivatives of
fixed order in a congruence-restricted twisting family \cite{Pomykala1997}.
In the family of weight-two newforms of varying prime level,
Kowalski, Michel, and VanderKam combined high completed derivatives,
mollification, and consecutive derivative orders to control both
functional-equation signs. Their passage from these nonvanishing counts
to rank-weighted tails also uses a separate bound for the second moment
of analytic rank
\cite[Sections~2, 4, and~8]{KowalskiMichelVanderKam2000}.
In the present quadratic-twist argument the derivative order grows with
height. This uniformity, together with a pointwise rank bound, controls
the rank mass beyond the available counting range without requiring a
bounded second analytic-rank moment.

Second moments of the twist \(L\)-functions themselves supply the
analytic comparison used below. For full-level holomorphic Hecke forms
of weight divisible by four, Soundararajan and Young obtained the
second-moment asymptotic over positive discriminants \(8d\), with \(d\)
odd and squarefree, under GRH for the twists, the Riemann zeta function,
and the symmetric-square \(L\)-function
\cite[Theorem~1.2]{SoundararajanYoung2010}.
Li proved the asymptotic unconditionally in this full-level setting
\cite[Theorem~1.1]{Li2022}. We adapt the Fourier-cutoff and prime-square
inflation mechanism in Li's proof, establishing the needed logarithmic
bounds for our finite family of arbitrary elliptic curves with its
conductor and bad-prime factors retained. Neither asymptotic is an input
to the argument.

\subsection{The analytic argument}

The starting point is the exact identity of Lemma~\ref{weights:vanishing},
forced by the vanishing of a derivative of order \(k\) of the completed \(L\)-function.
Its Dirichlet-series weight is obtained by exponentially smoothing
a \(k\)-th power of a truncated logarithm.
Near the square root of the conductor, at lengths of order \(X\) for
twists of height \(X\), this weight suppresses the late terms.
The approximation \((1-u)^k\approx e^{-ku}\) for \(u=O(1/k)\)
motivates mollification: at \(u=\log n/\log X\), the exponential
corresponds to a small positive shift from the center of the
\(L\)-function. A short approximation to its inverse Euler product
then makes the first weighted piece close to a positive main term.
We split the weighted series using smooth cutoffs into finite pieces
and a terminal series, use progressively shorter mollifiers as the
lengths of the finite pieces increase, and estimate the terminal series
without mollification.

Two features make the resulting estimates uniform for
\(k\) as large as a fixed power of \(\log X\).
First, a Poisson argument and inflation give unmollified second moments
with only powers of logarithms
(Propositions~\ref{analytic:moment} and~\ref{analytic:polynomial-moment}).
The inflation mechanism is related
to Li's work on quadratic twists
\cite[Lemmas~2.7 and~3.1, Proposition~3.2]{Li2022}; we prove the estimate
needed here for the entire fixed finite family of curves, including its
local factors at bad primes.
Second, Proposition~\ref{model:main-model-estimates} gives mean-square
bounds for the mollified finite pieces in a model that replaces
odd-prime character values by independent variables with their
complete-residue distributions. The mean-square estimates use the whole
probability space, without discarding exceptional outcomes.
Proposition~\ref{analytic:comparison} transfers these bounds to integer averages.
Separately, Proposition~\ref{model:good-event} compares the mollifiers
with positive reciprocal Euler products on most integer parameters
and relates these products to one another. We can then remove the
mollifiers and compare every piece with the same positive product:
the first piece dominates the later pieces and
the terminal series. The resulting nonzero sum contradicts the
completed-derivative identity whenever the rank exceeds the derivative
order of matching parity. Consecutive orders cover both signs.

For admissible twists of height \(X\), Proposition~\ref{tail:count} proves a bound
\(O_E(X/k^2)\) for the number of ranks exceeding \(k\), uniformly up to
\(k=(\log X)^{3/5}\).
A conductor-uniform Jensen bound (Lemma~\ref{tail:jensen}) gives the pointwise estimate
\(a(E^{(d)})=O_E(\log X)\).
The count at the largest permissible \(k\) then controls the remaining
rank mass, and summation of the integer tails proves
Theorem~\ref{thm:tail}.
The numerical exponents used in the length partition leave ample
room between the comparison error and the later thresholds; their
optimization is not needed.

The sections follow the inputs required by this argument.
Section~\ref{sec:setup} fixes the finite auxiliary family and the integer
residue model. Section~\ref{sec:analytic} proves the moments and the
comparison with that model. Section~\ref{sec:weights} constructs the
completed-derivative weights and bounds their terminal series.
Section~\ref{sec:model} constructs the short mollifiers and proves the
model estimates. Section~\ref{sec:tail} combines them into the counting
and rank-weighted tail bounds, and uses the companion density theorem
only in the final deduction of the mean. Section~\ref{sec:algebraic}
proves the separate algebraic-rank moment corollary.

\subsection{Conventions}

Implied constants may depend on the original curve \(E\), on fixed
smooth functions, and on a specified fixed order of differentiation.
They do not depend on the varying height, derivative order, or scale
index unless this is explicitly stated.
We write \(\tau(n)\) for the divisor function and use the Fourier kernel
\(e^{-2\pi ixy}\) when applying Poisson summation.
All sufficiently large lower thresholds for the integer \(k\) are
chosen after the fixed analytic exponents and smoothness orders.
